\documentclass[conference]{IEEEtran}
\IEEEoverridecommandlockouts
\usepackage{cite}
\usepackage{amsmath,amssymb}
\usepackage{graphicx}
\usepackage{textcomp}
\usepackage{xcolor}
\usepackage{booktabs}
\usepackage{url}
\def\BibTeX{{\rm B\kern-.05em{\sc i\kern-.025em b}\kern-.08em T\kern-.1667em\lower.7ex\hbox{E}\kern-.125emX}}

\begin{document}

\title{Design and Evaluation of a Five-Stage Pipelined RISC-V Processor
with Operand Forwarding and Dynamic Branch Prediction\\
{\footnotesize Modern Computer Architecture --- Phase I}}

% ===========================================================
% TODO before submission: replace the three placeholder names
% and all roll numbers with the real values.
% ===========================================================
\author{
\IEEEauthorblockN{Kunal Vats}
\IEEEauthorblockA{\textit{Dept. of Computer Applications} \\
\textit{{[Institution]}}\\
{[Roll No.]} \\
1kunalvats9@gmail.com}
\and
\IEEEauthorblockN{Dev Kothari}
\IEEEauthorblockA{\textit{Dept. of Computer Applications} \\
\textit{{[Institution]}}\\
{[Roll No.]} \\
devkothari32@gmail.com}
\and
\IEEEauthorblockN{{[Member 3]}}
\IEEEauthorblockA{\textit{Dept. of Computer Applications} \\
\textit{{[Institution]}}\\
{[Roll No.]} \\
{[email]}}
\and
\IEEEauthorblockN{{[Member 4]}}
\IEEEauthorblockA{\textit{Dept. of Computer Applications} \\
\textit{{[Institution]}}\\
{[Roll No.]} \\
{[email]}}
\and
\IEEEauthorblockN{{[Member 5]}}
\IEEEauthorblockA{\textit{Dept. of Computer Applications} \\
\textit{{[Institution]}}\\
{[Roll No.]} \\
{[email]}}
}

\maketitle

\begin{abstract}
This paper presents the design, implementation and quantitative evaluation
of a five-stage pipelined RISC-V processor implementing an RV32I subset
(\texttt{add}, \texttt{sub}, \texttt{addi}, \texttt{lw}, \texttt{sw},
\texttt{beq}, \texttt{bne}) in Verilog. Rather than building a single
processor, we construct \emph{four} cores from a common RTL base: a
single-cycle reference, and three pipelined configurations that differ
only by parameters---no hazard mitigation, operand forwarding, and
forwarding combined with a 64-entry two-bit dynamic branch predictor.
This parameterisation allows each architectural feature to be isolated and
measured independently. Correctness is established by a three-tier
methodology: unit tests, 34 directed hazard assertions, and a golden-model
comparison in which every pipelined configuration must reproduce the
single-cycle core's architectural state exactly. Across seven benchmarks
totalling 523 instructions, cycles per instruction fall from 2.302 for the
unmitigated pipeline to 1.620 with forwarding and 1.245 with prediction,
while branch prediction accuracy rises from 11.7\% to 93.3\%. Combining
cycle counts with gate-level clock-period estimates yields the paper's
central finding: relative to the single-cycle reference, the unmitigated
pipeline is \emph{slower} (0.71$\times$), forwarding achieves parity
(1.01$\times$), and only the addition of branch prediction produces a
genuine speedup (1.32$\times$). Pipelining alone is not a performance win;
pipelining together with hazard mitigation is.
\end{abstract}

\begin{IEEEkeywords}
RISC-V, instruction pipelining, data hazards, operand forwarding, branch
prediction, CPI, Verilog, processor verification
\end{IEEEkeywords}

% ===========================================================
\section{Introduction}

Instruction pipelining is the foundational technique of modern processor
design, and the five-stage pipeline---fetch, decode, execute, memory,
write-back---remains the standard pedagogical vehicle for teaching it
\cite{patterson2020}. Its appeal is that it decomposes a long
combinational path into stages that can operate concurrently, raising
throughput without increasing the work done per instruction.

The difficulty is that this concurrency is only correct if the dependencies
between overlapping instructions are respected. A pipelined datapath that
ignores hazards computes wrong answers; one that handles them by stalling
unconditionally computes right answers slowly. The engineering content of
the problem lies entirely in the mitigation mechanisms---operand
forwarding, hazard interlocks, and control-flow speculation---and in
determining whether they are worth their cost.

A common weakness in undergraduate and postgraduate implementations of
this exercise is that a single finished processor is built and a single
CPI figure is reported. Such a figure is difficult to interpret: it
confounds the cost of pipelining, the benefit of forwarding and the
benefit of prediction into one number, and it cannot answer whether any
individual mechanism justified its area.

\subsection{Contributions}

This work makes the following contributions:

\begin{enumerate}
\item A synthesisable five-stage pipelined RISC-V core supporting an RV32I
      subset, with full operand forwarding, load-use interlocking, and a
      two-bit dynamic branch predictor.
\item A \emph{four-configuration} methodology in which three pipelined
      variants are generated from identical RTL by parameter change alone,
      so that each measurement isolates exactly one architectural feature.
\item A three-tier verification strategy culminating in golden-model
      equivalence against a single-cycle reference, together with a
      cycle-accounting invariant that independently validates the
      performance counters.
\item A quantitative evaluation combining cycle counts with synthesised
      area and critical-path depth, including a sensitivity analysis that
      identifies the conditions under which pipelining becomes
      worthwhile.
\end{enumerate}

% ===========================================================
\section{Literature Review}

\subsection{The RISC-V Instruction Set}

RISC-V is an open, royalty-free instruction set architecture developed at
UC Berkeley and now maintained by RISC-V International
\cite{riscvisa}. Its base integer subset, RV32I, comprises a small,
regular, fixed-length instruction encoding designed explicitly to be
straightforward to implement in hardware. Register \texttt{x0} is
hardwired to zero, immediate fields are placed to maximise sign-bit
reuse, and there are no condition codes or implicit side effects. These
properties make RV32I well suited to a teaching implementation, and they
are the reason the decode and immediate-generation logic in this work is
comparatively simple.

Open-source implementations span a wide complexity range. PicoRV32
\cite{picorv32} is a size-optimised multi-cycle core intended for FPGA
soft-core use, while the Rocket Chip generator \cite{rocketchip} produces
cache-coherent out-of-order-capable designs. Our work sits at the
pedagogical end of this spectrum: the objective is measurement and
analysis of a classical pipeline, not competitive performance.

\subsection{Pipelining and Hazard Classification}

The canonical treatment of the five-stage pipeline and its hazards is due
to Patterson and Hennessy \cite{patterson2020}, with the quantitative
framework developed further in \cite{hennessy2017}. Hazards are classified
as structural (two instructions competing for one resource), data (an
instruction requiring a result not yet available), and control (the next
instruction address not yet known).

Among data hazards, only read-after-write (RAW) can occur in an in-order
single-issue pipeline; write-after-read and write-after-write arise only
when instructions complete out of order, as in the dynamically scheduled
designs introduced by Tomasulo \cite{tomasulo1967}. We confirm this
analysis for our design in Section III and use it to bound the
verification effort.

Operand forwarding---routing a result from a pipeline register directly to
a later instruction's execute input rather than waiting for write-back---is
the standard mitigation for RAW hazards. The literature is consistent that
one case resists it: a load followed immediately by a consumer of the
loaded value, where the data is simply not available in time
\cite{patterson2020}. This load-use hazard requires a one-cycle
interlock, and our measurements in Section V confirm it as the dominant
residual overhead once all other mechanisms are in place.

\subsection{Branch Prediction}

Control hazards are addressed by predicting branch outcomes. Smith's
foundational study \cite{smith1981} evaluated several strategies and
established the two-bit saturating counter as an effective and inexpensive
dynamic predictor. The key insight is hysteresis: a single-bit predictor
mispredicts twice per loop, once on the terminating iteration and again on
the first iteration of the next execution, because the terminating
iteration has already flipped the bit. A two-bit counter absorbs one
anomalous outcome without changing its prediction.

Subsequent work increased accuracy substantially. Yeh and Patt
\cite{yehpatt1991} introduced two-level adaptive predictors that correlate
on branch history, and McFarling \cite{mcfarling1993} showed that combining
predictors with a selector outperforms either component. These schemes
exploit correlation between branches, which the loop-dominated benchmarks
used in this work do not exhibit; we therefore implement the two-bit
scheme of \cite{smith1981} and discuss the correlating alternatives as
future work.

\subsection{Verification Methodology}

Processor verification by comparison against a reference model is
long-established industrial practice, and is the basis of the RISC-V
architectural test framework \cite{riscvarchtest}. The essential idea is
that the design under test and a trusted reference execute identical
programs, after which their architecturally visible state is compared. We
adopt this approach, using our own single-cycle implementation as the
reference, which has the advantage that the two share an instruction
decoder and can therefore be compared instruction-for-instruction without
an ISA simulator.

\subsection{Positioning of This Work}

The individual mechanisms described above are well understood and
textbook-standard. The contribution here is methodological: by generating
all pipelined configurations from one parameterised source, the marginal
cost and marginal benefit of each mechanism can be attributed
unambiguously. We are not aware of this being standard practice in
teaching implementations, and Section V shows it producing a conclusion
that a single-configuration study would miss entirely.

% ===========================================================
\section{Proposed Solution}

\subsection{Instruction Set Subset}

We implement \texttt{add}, \texttt{sub}, \texttt{addi}, \texttt{lw},
\texttt{sw}, \texttt{beq} and \texttt{bne}, with \texttt{nop} encoded
canonically as \texttt{addi x0, x0, 0}. This subset exercises every
hazard class---register-to-register dependencies, memory dependencies and
control flow---while keeping the decoder small enough that it is not a
source of defects. Jumps, shifts, logical operations and comparisons are
excluded; their consequences are discussed in Section VI.

\subsection{Microarchitecture}

The datapath is partitioned into the five classical stages separated by
four pipeline registers (IF/ID, ID/EX, EX/MEM, MEM/WB). Several design
decisions are worth recording.

\textbf{Harvard memory organisation.} Instruction and data memories are
separate, which eliminates the structural hazard between fetch and memory
access entirely. A unified memory would force a stall on every load and
store.

\textbf{Control signals travel with data.} All control is decoded once, in
ID, and carried down the pipeline inside the pipeline registers alongside
the operands. Control is never recomputed in a later stage. Recomputing it
from whichever instruction a later stage happens to hold is a classical
defect in which control is applied to the wrong instruction.

\textbf{Register file with write-to-read bypass.} The register file
provides two combinational read ports and one synchronous write port, with
\texttt{x0} hardwired to zero. An internal bypass returns the value being
written in the same cycle, which removes the WB-to-ID hazard before the
hazard logic observes it, reducing the forwarding network to two sources.
This bypass is parameterised, because in the single-cycle core the read and
write belong to the \emph{same} instruction and the bypass would close a
combinational loop.

\textbf{Branch resolution in EX.} Resolving conditional branches in the
execute stage rather than in memory reduces the misprediction penalty from
three wasted fetch slots to two.

\subsection{Hazard Analysis}

For this in-order, single-issue, single-write-port design we establish that
write-after-read and write-after-write hazards cannot occur: instructions
issue in order, each writes at most one register, and all writes occur in
write-back in program order. Structural hazards cannot occur given separate
memories and a register file with exactly the required port count. The
design must therefore handle RAW data hazards and control hazards, and
nothing else.

\subsection{Operand Forwarding}

Two bypass paths are provided: EX/MEM to EX for a distance-one dependency,
and MEM/WB to EX for distance two. Dependencies at distance three and
beyond are satisfied by the register file bypass described above.

Two correctness conditions are essential. First, when both pipeline
registers hold writes to the same architectural register, the EX/MEM value
must take priority, since it belongs to the younger instruction whose write
occurs later in program order. Second, forwarding from a destination of
\texttt{x0} must be suppressed, because such writes are discarded and the
forwarded value would never become architecturally visible.

\subsection{Load-Use Interlock}

A load in EX whose destination is read by the instruction in ID constitutes
the one RAW hazard forwarding cannot resolve: the loaded value is returned
\emph{during} the cycle in which the consumer requires it at its input. One
bubble is therefore unavoidable, after which the MEM/WB path supplies the
value. Detection is qualified by signals indicating whether the instruction
in ID genuinely reads each source register, since for \texttt{addi} and
\texttt{lw} the \texttt{rs2} field is part of the immediate; without this
qualification the interlock would stall on false dependencies and silently
inflate CPI.

\subsection{Dynamic Branch Prediction}

The predictor is a 64-entry table of two-bit saturating counters indexed by
\texttt{PC[7:2]}, read in IF and updated in EX when a branch resolves,
following \cite{smith1981}. Counters reset to the weakly-not-taken state,
so the first prediction for any branch matches what a static not-taken
policy would produce; the comparison between configurations therefore
isolates the effect of learning.

No branch target buffer is required. Because the instruction memory read is
combinational, the instruction word is available in IF, so the opcode
identifies a branch and a B-format immediate decoder with one adder
produces the target directly. The cost is additional logic depth in IF,
quantified in Section V.

Prediction is verified in EX by comparing the recorded prediction against
the computed outcome. A predictor capable of predicting \emph{taken} can be
wrong in either direction, so the redirect target is a multiplexer
selecting the branch target or the fall-through address, rather than the
branch target alone as static not-taken prediction would require.

\subsection{Flush Safety}

When a bubble is inserted, every control bit is cleared, not only the
validity bit. Were validity alone cleared, a flushed store would retain an
asserted memory-write enable and would corrupt data memory---a defect
whose symptom appears only much later, when the corrupted location is read
back. A directed test exists specifically to detect this condition.

\subsection{The Four-Configuration Methodology}

Table~\ref{tab:configs} summarises the configurations. The three pipelined
variants are produced from identical source by parameter override, so any
difference in measured cycles or synthesised area is attributable to the
single feature that changed.

\begin{table}[htbp]
\caption{Evaluated Configurations}
\begin{center}
\begin{tabular}{llp{3.6cm}}
\toprule
\textbf{Name} & \textbf{Parameters} & \textbf{Role} \\
\midrule
\texttt{single} & --- & Single-cycle golden reference; CPI $=1.000$ by construction \\
\texttt{stall}  & \texttt{FORWARD\_EN=0} \texttt{PREDICTOR=0} & Pipeline with no hazard mitigation \\
\texttt{fwd}    & \texttt{FORWARD\_EN=1} \texttt{PREDICTOR=0} & Adds the forwarding network \\
\texttt{pred}   & \texttt{FORWARD\_EN=1} \texttt{PREDICTOR=1} & Adds the two-bit predictor \\
\bottomrule
\end{tabular}
\label{tab:configs}
\end{center}
\end{table}

% ===========================================================
\section{Implementation and Verification}

\subsection{Toolchain}

The design is written in Verilog-2001 and simulated with Icarus Verilog
\cite{iverilog}; waveforms are inspected in GTKWave. Technology-independent
synthesis uses the Yosys Open SYnthesis Suite \cite{yosys} to obtain cell
counts and longest-topological-path figures. An FPGA flow targeting a
Xilinx Artix-7 device is prepared for Phase II.

A purpose-written assembler translates the supported subset into hex memory
images, emitting a listing that pairs each encoding with its source line.
Programs terminate by a simulation-only convention in which the word
\texttt{0x00000000}---not a legal RISC-V encoding---halts fetch, allowing
retired instructions to be counted exactly without prior knowledge of
program length.

\subsection{Three-Tier Verification}

\textbf{Unit tests} cover the ALU and the register file, the latter
specifically confirming that \texttt{x0} reads as zero, discards writes,
and does not participate in the write-to-read bypass.

\textbf{Directed hazard tests} comprise 34 assertions grouped by hazard
class: distance-one and distance-two forwarding, forwarding priority,
\texttt{x0} exclusion, the load-use interlock, store-data forwarding,
not-taken and taken branches, operand forwarding into the branch
comparator, and the flushed-store case described in Section III-G. Because
the expected values are architectural, the same assertions pass on all four
configurations irrespective of how many bubbles each requires.

\textbf{Golden-model comparison} is the decisive check. Each configuration
executes every benchmark and emits an architectural signature comprising
registers \texttt{x1}--\texttt{x31} and the first 32 words of data memory.
Every pipelined signature must match the single-cycle signature exactly.

\subsection{Cycle-Accounting Invariant}

For every pipelined run the harness verifies
\begin{equation}
C = I + 4 + S + 2M
\label{eq:accounting}
\end{equation}
where $C$ is elapsed cycles, $I$ retired instructions, $4$ the pipeline
fill, $S$ stall cycles and $M$ mispredictions. This holds exactly on all 21
pipelined runs. It is a strong independent check: an unaccounted cycle
would indicate a stall or flush taken but not reported, and a missing one
would indicate a defective counter.

\subsection{Defects Found}

Three defects were found by this infrastructure rather than by inspection,
and are recorded because verification quality is itself an evaluation
criterion. First, the register-file bypass closed a combinational loop in
the single-cycle core for any instruction reusing its destination as a
source; three benchmarks passed before one triggered it. Second, a
testbench race in reset deassertion made cycle counts non-deterministic by
$\pm1$, sufficient to corrupt every CPI figure. Third, under
\texttt{FORWARD\_EN=0} the forwarding unit's sensitivity list constant-folded
to empty, leaving its outputs undefined throughout simulation; results
remained correct only because an undefined case selector falls through to a
default branch that happened to be correct.

% ===========================================================
\section{Results}

All measurements are over seven benchmarks totalling 523 instructions,
spanning dependency-free code, dependency chains, load-use sequences,
counted loops and branch-dominated control flow.

\subsection{Cycles per Instruction}

\begin{table}[htbp]
\caption{CPI and IPC by Configuration (523 instructions)}
\begin{center}
\begin{tabular}{lrrrr}
\toprule
\textbf{Configuration} & \textbf{Cycles} & \textbf{CPI} & \textbf{IPC} & \textbf{Saved} \\
\midrule
Single-cycle          & 523  & 1.000 & 1.000 & --- \\
Pipelined, stall-only & 1204 & 2.302 & 0.434 & --- \\
Pipelined, forwarding & 847  & 1.620 & 0.617 & 357 \\
Pipelined, fwd + pred & 651  & 1.245 & 0.803 & 196 \\
\bottomrule
\end{tabular}
\label{tab:cpi}
\end{center}
\end{table}

Table~\ref{tab:cpi} shows that the two mitigation mechanisms together
remove 553 of the 681 hazard cycles incurred by the unmitigated pipeline,
or 81\% of the overhead. On the longest benchmark, a 64-element vector sum
of 260 instructions, steady-state CPI is 1.277; on an iterative Fibonacci
computation it reaches 1.076, within 7\% of the ideal for a scalar
pipeline.

\subsection{Branch Prediction Accuracy}

\begin{table}[htbp]
\caption{Branch Prediction Accuracy}
\begin{center}
\begin{tabular}{lrrrr}
\toprule
\textbf{Benchmark} & \textbf{Br.} & \textbf{Taken} & \textbf{Static} & \textbf{2-bit} \\
\midrule
\texttt{array\_sum} & 16  & 15  & 6.2\%  & 87.5\% \\
\texttt{branch\_mix}& 20  & 9   & 55.0\% & 90.0\% \\
\texttt{fib}        & 20  & 19  & 5.0\%  & 90.0\% \\
\texttt{vecsum64}   & 64  & 63  & 1.6\%  & 96.9\% \\
\midrule
\textbf{Total}      & 120 & 106 & 11.7\% & \textbf{93.3\%} \\
\bottomrule
\end{tabular}
\label{tab:bp}
\end{center}
\end{table}

Accuracy rises from 11.7\% to 93.3\%, eliminating 98 of 106 mispredictions
and 196 cycles of flushed work. Notably, the predictor mispredicts
\emph{exactly twice per loop} on every benchmark: once on entry, because
the counter initialises to weakly-not-taken, and once on exit. This matches
the behaviour predicted by \cite{smith1981} precisely.

\subsection{Synthesised Area}

\begin{table}[htbp]
\caption{Marginal Cost of Each Mechanism (Yosys, memories excluded)}
\begin{center}
\begin{tabular}{lrr}
\toprule
\textbf{Change} & \textbf{Cells} & \textbf{Flip-flops} \\
\midrule
Single-cycle $\rightarrow$ pipelined & $+1497$ & $+413$ \\
Adding operand forwarding            & $+288$  & $\mathbf{+0}$ \\
Adding the two-bit predictor         & $+1071$ & $+136$ \\
\bottomrule
\end{tabular}
\label{tab:area}
\end{center}
\end{table}

Operand forwarding is the most favourable mechanism in the design: 288
gates and \emph{no} sequential state, because it is purely combinational,
in exchange for a 1.42$\times$ cycle-count improvement. The predictor's 136
flip-flops are almost entirely the counter table itself (64 entries
$\times$ 2 bits), with two further bits carrying the prediction down the
pipeline.

Per-stage path analysis identifies EX as the deepest stage at 70 gate
levels, of which the 32-bit ALU adder alone accounts for 67. The predictor
increases IF depth from 30 to 56 levels, which remains below EX and
therefore does not lengthen the critical path in this implementation.

\subsection{The Central Result}

Cycle counts alone understate pipelining, because the single-cycle core
requires a clock period covering all five stages in series whereas the
pipelined core requires only the deepest single stage. The appropriate
comparison is
\begin{equation}
\text{speedup} = \frac{C_{\text{single}}}{C_{\text{cfg}}}
                 \times \frac{T_{\text{single}}}{T_{\text{cfg}}}
\label{eq:speedup}
\end{equation}

\begin{table}[htbp]
\caption{Net Speedup Relative to the Single-Cycle Reference}
\begin{center}
\begin{tabular}{lrrr}
\toprule
\textbf{Configuration} & \textbf{Cycle} & \textbf{Clock} & \textbf{Net} \\
\midrule
Pipelined, stall-only & 0.434$\times$ & 1.64$\times$ & 0.71$\times$ \\
Pipelined, forwarding & 0.617$\times$ & 1.64$\times$ & 1.01$\times$ \\
Pipelined, fwd + pred & 0.803$\times$ & 1.64$\times$ & \textbf{1.32}$\times$ \\
\bottomrule
\end{tabular}
\label{tab:speedup}
\end{center}
\end{table}

Table~\ref{tab:speedup} is the principal result of this work. Read in
sequence, the unmitigated pipeline is \emph{slower} than the single-cycle
reference it replaced; forwarding restores parity; and only with branch
prediction does the design become faster. \textbf{Pipelining is not by
itself a performance improvement---pipelining together with hazard
mitigation is.} A study reporting only the final configuration would
observe a 1.32$\times$ speedup and attribute it to pipelining, concealing
the fact that the mitigation mechanisms, not the pipeline, are responsible
for the entire gain.

\subsection{Sensitivity to Memory Latency}

Because both cores are evaluated with memories excluded from the timing
path, the single-cycle core is flattered: its critical path contains an
instruction fetch and a data access in series, whereas no single pipeline
stage contains both. Modelling a memory read delay $M$ in gate levels, the
net speedup rises from 1.32$\times$ at $M=0$ to 1.64$\times$ at $M=20$,
beyond which the critical stage migrates from EX to IF and the ratio
plateaus. This is a quantitative statement of why practical pipelines
require caches and why deeper pipelines subdivide fetch and memory access.

% ===========================================================
\section{Individual Contributions}

% ===========================================================
% TODO: replace placeholder names. If the team has fewer than
% five members, merge rows as follows: Front End + Decode,
% Execute + Back End, Verification + Infrastructure.
% ===========================================================

\begin{table}[htbp]
\caption{Division of Work}
\begin{center}
\begin{tabular}{p{1.9cm}p{2.1cm}p{3.3cm}}
\toprule
\textbf{Member} & \textbf{Area} & \textbf{Modules and artefacts} \\
\midrule
{[Member 3]} & Front end &
\texttt{pc\_if.v}, \texttt{if\_id.v}, \texttt{imem.v},
\texttt{branch\_predictor.v}; next-PC selection, prediction lookup,
flush handling \\
\addlinespace
Dev Kothari & Decode &
\texttt{decode\_stage.v}, \texttt{id\_ex.v}, \texttt{regfile.v},
\texttt{imm\_gen.v}, \texttt{control\_unit.v}; control bundle,
immediate formats, register-file bypass \\
\addlinespace
{[Member 4]} & Execute &
\texttt{ex\_stage.v}, \texttt{alu.v}, \texttt{branch\_cmp.v},
\texttt{forwarding\_unit.v}, \texttt{ex\_mem.v}; bypass network,
prediction check and redirect \\
\addlinespace
Kunal Vats & Back end &
\texttt{mem\_stage.v}, \texttt{mem\_wb.v}, \texttt{wb\_stage.v},
\texttt{dmem.v}, \texttt{hazard\_unit.v}; load-use interlock, stall and
bubble generation, flush safety \\
\addlinespace
{[Member 5]} & Verification and infrastructure &
\texttt{tb\_*.v}, assembler, benchmark programs, regression and synthesis
scripts, results tooling \\
\bottomrule
\end{tabular}
\label{tab:contrib}
\end{center}
\end{table}

Every RTL source file records its owner in a header comment, and the
division above is reflected in the repository commit history.

% ===========================================================
\section{Limitations and Future Work}

The instruction subset excludes jumps, and therefore function calls;
large-constant formation via \texttt{lui}; shifts and logical operations;
and magnitude comparison, which is why no sorting benchmark is included.
There are no exceptions, interrupts or privilege modes, and both memories
are idealised single-cycle structures without caching.

The data memory uses an asynchronous read, which is required for the
loaded value to be captured by MEM/WB within a single memory stage but
prevents inference of FPGA block RAM. The predictor is direction-only and
untagged, so widely separated branches may alias; no benchmark here is
large enough to exhibit this.

Three directions follow naturally. First, FPGA implementation to obtain
measured $F_{\max}$ in place of the gate-level clock-period estimate used
in Table~\ref{tab:speedup}; this will also determine whether the predictor
remains free in clock period once adders map onto dedicated carry chains.
Second, a correlating predictor in the style of \cite{yehpatt1991} or
\cite{mcfarling1993}. Third, and indicated most strongly by our own
results, compiler scheduling to fill the load-use delay slot, which is now
the dominant remaining source of lost cycles.

% ===========================================================
\section{Conclusion}

We have designed, verified and evaluated a five-stage pipelined RISC-V
processor in four configurations generated from a single parameterised
source. Verification combines unit tests, 34 directed hazard assertions and
golden-model equivalence against a single-cycle reference, with an
independent cycle-accounting invariant holding exactly across all runs.
Operand forwarding and two-bit dynamic branch prediction together remove
81\% of hazard overhead, reducing CPI from 2.302 to 1.245 and raising
branch prediction accuracy from 11.7\% to 93.3\%.

The principal finding is that these mechanisms are not refinements but
prerequisites: without them the pipelined design is measurably slower than
the single-cycle core it replaces. Isolating each mechanism by construction
was what made this visible, and we regard the methodology as the more
transferable of the two contributions.

% ===========================================================
\begin{thebibliography}{00}

\bibitem{riscvisa} A. Waterman and K. Asanovi\'{c}, Eds., \emph{The RISC-V
Instruction Set Manual, Volume I: Unprivileged ISA}, RISC-V
International, 2019.

\bibitem{patterson2020} D. A. Patterson and J. L. Hennessy, \emph{Computer
Organization and Design RISC-V Edition: The Hardware/Software Interface},
2nd ed. Cambridge, MA, USA: Morgan Kaufmann, 2020.

\bibitem{hennessy2017} J. L. Hennessy and D. A. Patterson, \emph{Computer
Architecture: A Quantitative Approach}, 6th ed. Cambridge, MA, USA: Morgan
Kaufmann, 2017.

\bibitem{smith1981} J. E. Smith, ``A study of branch prediction
strategies,'' in \emph{Proc. 8th Annu. Symp. Computer Architecture (ISCA)},
1981, pp. 135--148.

\bibitem{yehpatt1991} T.-Y. Yeh and Y. N. Patt, ``Two-level adaptive
training branch prediction,'' in \emph{Proc. 24th Annu. Int. Symp.
Microarchitecture (MICRO)}, 1991, pp. 51--61.

\bibitem{mcfarling1993} S. McFarling, ``Combining branch predictors,''
Digital Equipment Corp. Western Research Laboratory, Tech. Note TN-36,
1993.

\bibitem{tomasulo1967} R. M. Tomasulo, ``An efficient algorithm for
exploiting multiple arithmetic units,'' \emph{IBM Journal of Research and
Development}, vol. 11, no. 1, pp. 25--33, 1967.

\bibitem{rocketchip} K. Asanovi\'{c} \emph{et al.}, ``The Rocket Chip
generator,'' EECS Dept., Univ. California, Berkeley, Tech. Rep.
UCB/EECS-2016-17, 2016.

\bibitem{picorv32} C. Wolf, ``PicoRV32: A size-optimized RISC-V CPU.''
[Online]. Available: \url{https://github.com/YosysHQ/picorv32}

\bibitem{yosys} C. Wolf, ``Yosys Open SYnthesis Suite.'' [Online].
Available: \url{https://yosyshq.net/yosys/}

\bibitem{iverilog} S. Williams, ``Icarus Verilog.'' [Online]. Available:
\url{http://iverilog.icarus.com/}

\bibitem{riscvarchtest} RISC-V International, ``RISC-V Architectural Test
Suite.'' [Online]. Available:
\url{https://github.com/riscv-non-isa/riscv-arch-test}

\end{thebibliography}

\end{document}
